\documentclass{report}
\usepackage[utf8]{inputenc}
\usepackage{polski}
\usepackage[T1]{fontenc}
\usepackage{graphicx}
\usepackage{amsmath}
\usepackage{amssymb}

\title{Święty Mikołaj}
\author{Jakub Ronkiewicz}
\begin{document}
\maketitle

\chapter{Historia postaci}
\section{Ogólne informacje}
\includegraphics[width=12cm,keepaspectratio]{1.png}\newline
Święty Mikołaj – postać starszego mężczyzny z białą brodą ubranego w czerwony strój, który wedle różnych legend i baśni w okresie świąt Bożego Narodzenia rozwozi dzieciom prezenty saniami ciągniętymi przez zaprzęg reniferów. Według różnych wersji zamieszkuje wraz z grupą elfów Laponię lub biegun północny.Dzięki sprawnej promocji praktycznie zastąpił on w powszechnej świadomości tradycyjny wizerunek świętego Mikołaja, biskupa Miry.Obecnie powszechna forma tej postaci wywodzi się z kultury brytyjskiej i amerykańskiej, gdzie jest jedną z atrakcji bożonarodzeniowych. W większej części Europy, w tym w Polsce, Mikołajki obchodzone są tradycyjnie 6 grudnia jako wspomnienie świętego Mikołaja, biskupa Miry. Rankiem tego dnia dzieci, które przez cały mijający rok były grzeczne, znajdują prezenty, ukryte pod poduszką, w buciku lub w innym specjalnie przygotowanym w tym celu miejscu (np. w skarpecie).\newpage
\section{Lokalne odpowiedniki}
W Rosji św. Mikołaj był również otoczony kultem jako patron uciśnionych i skrzywdzonych, ale przejętą z zachodu postać rozdającą dzieciom prezenty nazwano Dziadkiem Mrozem, ze względu na podobieństwo do bajkowej postaci nazywanej Moroz Krasnyj Nos. Zgodnie z tradycją, rozdaje on prezenty wraz ze Śnieżynką. W tradycji bizantyńskiej jego odpowiednikiem jest Święty Bazyli, który obdarowuje prezentami dzieci w dniu 1 stycznia.W Wielkopolsce, na Kujawach, na Pałukach, na Kaszubach i na Pomorzu Zachodnim prezenty na Boże Narodzenie tradycyjnie przynosi Gwiazdor, święty Mikołaj przynosi prezenty 6 grudnia.W Małopolsce prezenty pod choinką zostawia Aniołek, natomiast w nocy z 5 na 6 grudnia prezenty przynosi święty Mikołaj. Na Górnym Śląsku i w Czechach prezenty na Boże Narodzenie przynosi Dzieciątko utożsamiane z postacią Jezusa Chrystusa. 6 grudnia prezenty przynosi święty Mikołaj.Na Dolnym Śląsku oraz na Opolszczyźnie prezenty w Boże Narodzenie przynosi Gwiazdka. Mieszkańcom Płaskowyżu Tarnogrodzkiego natomiast na Boże Narodzenie prezenty przynoszą krasnoludki. Mikołaj jedynie 6 grudnia.
\chapter{Wzory matematyczne}
\begin{enumerate}
\item $f(x)=ax^2+bx+c,$ gdzie $a \ne 0$
\item $\binom{n}{k}=\frac{n!}{k\cdot(n-k)!}$
\item
\item
\item
\end{enumerate}

\end{document}
